%!TEX root = ../../main.tex
\section{선정, 보정, 동결}
\label{sec:pred-selection}

후보와 기준선은 15.14 TRAIN(out-of-fold 라벨)에서 적합되고, 15.15 Q\_SELECT에서 전체 T의 경기 가중 Brier로 선정된 뒤 동결되었다. Q\_CAL은 LightGBM 후보의 조기 종료에만 쓰였다. 모든 모델에 대해 Q\_CAL에서 양의 기울기 시그모이드 $g_q$를 적합하여 Q\_SELECT에서 항등과 비교하였고, $q$와 $\PTflex$에 대해서는 항등이 유지되었다. 같은 단계가 S 역할에서 따로 수행되었다(학습기는 \texttt{logit\_state}로 사전 고정, $\PTflexS$의 매듭과 $C$는 S Q\_SELECT에서 선정, 보정기는 계약에 따라 항등). 어떤 후보도 TEST에서 선정되지 않았다.

본 논문의 주 대비는 코호트당 하나씩이다. 같은 T 행에서 동결 $q$ 대 $\PTflex$의 $\dBrier$, 같은 S 행에서 $\qS$ 대 $\PTflexS$의 $\dBrier$이다. $\PTlin$은 이전 표와의 연속성을 위해 유지하고, 상수와 $p$-only 스플라인은 맥락으로만 보고한다. Transformer, 그래프 신경망, TabM 등으로 후보를 넓히지 않았다. 이는 실패에서 비롯된 결정이 아니라 선정 과적합의 위험과 동결의 원칙에서 비롯된 결정이다.
